\documentclass[11pt,a4paper]{article}

\usepackage[margin=1in]{geometry}
\usepackage{amsmath,amssymb,amsthm}
\usepackage{graphicx}
\usepackage{booktabs}
\usepackage{microtype}
\usepackage[numbers,sort&compress]{natbib}
\usepackage[colorlinks=true,linkcolor=blue,citecolor=blue,urlcolor=blue]{hyperref}

\newcommand{\ket}[1]{\lvert #1\rangle}
\newcommand{\bra}[1]{\langle #1\rvert}
\newcommand{\braket}[2]{\langle #1\vert #2\rangle}
\newcommand{\bramat}[3]{\langle #1\vert #2\vert #3\rangle}
\newcommand{\Tr}{\operatorname{Tr}}
\newcommand{\detmat}{\operatorname{det}}

\title{Matrix-Free Deflation for Frozen-Subspace\\ Neural Excited-State Variational Monte Carlo}

% TODO: add authors and affiliations before submission.
\author{}
\date{\today}

\begin{document}

\maketitle

\begin{abstract}
Neural-network quantum states have turned variational Monte Carlo (VMC) into a competitive electronic-structure method, but extending it to many electronic excited states remains expensive. State-specific orthogonalization requires explicit overlap information, while natural excited-state VMC (NES-VMC) optimizes $K$ states jointly through a determinant of $K$ replicated network outputs, so sampling, local-energy, and gradient costs grow with $K$ and already converged low-lying states continue to be perturbed when higher states are added. We introduce matrix-free deflation for frozen-subspace NES-VMC. Converged low-lying states are frozen as a projection basis, and a small number of new states are optimized in the orthogonal complement, so that adding high-lying states does not disturb the accurate low-lying ones. The central construction removes the need for full Hilbert-space vectors or a dense Hamiltonian: the projection coefficients are obtained from Monte Carlo estimates of the Gram matrices between frozen and active columns, accumulated on an auxiliary Markov-chain distribution with importance weights, where a common normalization constant cancels. The projected columns enter a small determinant and local-energy matrix, and automatic differentiation of the determinant automatically annihilates the frozen directions, yielding the projected gradient without constructing a projection operator. Final energies and eigenstates are extracted from a generalized eigenvalue problem estimated on consistently weighted samples. On H$_2$/6-31G (eight spin orbitals; 16-dimensional FCI space), freezing four states and optimizing two additional states recovers $E_4$ and $E_5$ to within $0.013$ and $0.025$ mHa with FCI overlaps of $0.99995$ and $0.99751$; optimizing four additional states recovers $E_4$--$E_7$ to within $0.003$--$0.010$ mHa. A direct $K=6$ optimization with a 300-step budget instead left the errors in $E_4$ and $E_5$ at $59.7$ and $493$ mHa. The training loop never forms a $2^N$-dimensional vector or Hamiltonian, and its cost grows with the number of active states rather than with Hilbert-space dimension. We also show that eigenvalues of the sample-averaged local-energy matrix are not reliable subspace-energy estimators for $K\ge 2$, and that the consistently weighted Gram generalized-eigenvalue estimator is required. These results establish frozen-subspace recycling as a practical route to systematically improvable excited-state calculations with neural quantum states.
\end{abstract}

\noindent\textbf{Keywords:} neural quantum states, variational Monte Carlo, excited states, deflation, matrix-free methods, natural excited states

\section{Introduction}

The accurate determination of electronic excited states is a central task in quantum chemistry, spectroscopy, photochemistry, and materials modeling. Full configuration interaction (FCI) provides the exact solution within a given one-electron basis \citep{knowles1984}, but its factorial cost restricts it to small active spaces. Quantum Monte Carlo (QMC) methods avoid explicit high-dimensional linear algebra and offer a favorable scaling with system size \citep{foulkes2001,becca2017}; variational Monte Carlo (VMC) in particular provides an explicit, systematically improvable wave function and an energy that is variational for the ground state.

The accuracy of VMC is limited by the flexibility of its trial wave function. The introduction of neural-network quantum states (NQS) \citep{carleo2017} removed much of this limitation: modern neural-network wave functions can represent many-body correlations compactly and have reached chemical accuracy for small molecules and solids. In an NQS-VMC calculation, the wave function is evaluated pointwise on configurations generated by Markov-chain Monte Carlo (MCMC), the Hamiltonian is applied through its sparse connection graph, and the parameters are optimized by stochastic gradient descent or stochastic reconfiguration (natural gradient) \citep{becca2017,vicentini2022}. Neither the wave function nor the Hamiltonian is ever stored as a $2^N$-dimensional object, which is the essential reason why the approach can address Hilbert spaces far beyond exact diagonalization.

Excited states are substantially more difficult. A variational energy alone does not select a specific excited state, because the Rayleigh quotient is minimized by the ground state. Common remedies include orthogonalization against lower states, state-averaged or penalty-based functionals, and symmetry- or variance-based objectives. Orthogonalization requires overlap information between the target state and all lower states, while penalty methods add a positive multiple of the squared overlap with lower states to the energy \citep{zhao2020}. Neural-network implementations have exploited symmetry sectors \citep{choo2018}, variance-based or penalty-based objectives \citep{entwistle2023}, and, more recently, natural excited states \citep{pfau2024}.

Natural excited-state VMC (NES-VMC) \citep{pfau2024} organizes the calculation around $K$ single-state wave functions $\psi_0,\ldots,\psi_{K-1}$; in the formulation used here, they are combined into a determinant on a $K$-fold replicated configuration space,
\begin{equation}
\Psi_\theta(\mathbf{x}_1,\ldots,\mathbf{x}_K)
= \det\bigl[\psi_j(\mathbf{x}_i)\bigr]_{i,j=0}^{K-1}.
\label{eq:nes}
\end{equation}
For each replica configuration the amplitudes and the Hamiltonian action form the $K\times K$ matrices
\begin{equation}
A_{ij}=\psi_j(\mathbf{x}_i),
\qquad
Z_{ij}=(H\psi_j)(\mathbf{x}_i),
\qquad
E_L=A^{-1}Z,
\label{eq:local}
\end{equation}
and the training loss is
\begin{equation}
\mathcal{L}(\theta)=\operatorname{Re}\Tr\langle E_L\rangle .
\label{eq:loss}
\end{equation}
At convergence, physical states and energies are obtained from the subspace Ritz problem
\begin{equation}
M v=\lambda S v,
\qquad
S_{ij}=\braket{\psi_i}{\psi_j},
\qquad
M_{ij}=\bramat{\psi_i}{H}{\psi_j},
\label{eq:ritz}
\end{equation}
so that $K$ eigenstates are approximated simultaneously. The determinant enforces linear independence of the columns and permits a single joint MCMC run.

In practice, NES-VMC becomes increasingly cumbersome as $K$ grows. Each additional state enlarges the replicated walker, the local-energy matrix, the number of trainable parameters, and the dimension of the stochastic-reconfiguration linear system. More importantly, the states converge at different rates. Once the low-lying states are accurate, continuing to optimize all $K$ states jointly exposes them to stochastic updates driven by the still-unconverged high-lying states; near-degeneracies and ill-conditioned determinants amplify this coupling. In our H$_2$/6-31G reference calculation, a direct $K=6$ optimization with a 300-step budget left the errors in $E_4$ and $E_5$ at $59.7$ and $493$ mHa, respectively, even though the same machinery had already produced an accurate $K=4$ solution. Adding a state should not require re-optimizing, and thereby risking, the states that are already correct.

This is a \emph{subspace recycling} problem. In numerical linear algebra and quantum chemistry, the standard tool is deflation: once a set of eigenpairs is known, one either projects the operator onto the orthogonal complement of the known invariant subspace or shifts the known eigenvalues away \citep{saad2011,wang1994}. Deflation is routinely used in iterative eigensolvers and in embedded correlated methods precisely because it allows new eigenpairs to be computed without recomputing the old ones. Transferring the idea to NQS-VMC, however, is not straightforward. First, a neural-network state is not stored as a vector; it exists only as a function that can be evaluated pointwise. Second, the projection coefficients are full Hilbert-space inner products, $\braket{\phi_i}{\psi_m}$, whose direct evaluation would reintroduce the exponential cost that NQS avoid. Third, the Hamiltonian is matrix-free and available only through local connections, so a projected Hamiltonian cannot be applied as a dense operator. Fourth, the sampling distribution, the local-energy estimator, and the parameter gradient must all remain mutually consistent after projection. Fifth, because the frozen states are only approximate eigenstates, the projected problem has an intrinsic bias relative to the exact eigenpairs, which must be separated from optimization error.

This paper develops a matrix-free deflation framework for frozen-subspace NES-VMC that addresses these issues. The main idea is summarized in Eqs.~\eqref{eq:projection}--\eqref{eq:extraction}. Suppose that $r$ low-lying states $\phi_0,\ldots,\phi_{r-1}$ have been converged and frozen. They are represented not by $2^N$-dimensional vectors but by the old network parameters together with a small rotation matrix, so that $\phi_i(x)$ can be evaluated on any sampled configuration $x$. New active columns $\psi_m$ are optimized only through their components in the orthogonal complement of the frozen span,
\begin{equation}
\chi_m=\psi_m-\sum_{i=0}^{r-1}c_{im}\phi_i,
\qquad
c=G^{-1}M,
\qquad
G_{ij}=\braket{\phi_i}{\phi_j},
\quad
M_{im}=\braket{\phi_i}{\psi_m}.
\label{eq:projection}
\end{equation}
The Gram matrices $G$ and $M$ are estimated on an auxiliary MCMC distribution
$q(y)\propto\sum_i|\phi_i(y)|^2+\varepsilon\sum_m|\psi_m(y)|^2$ with importance weights $1/q(y)$. All Gram blocks carry the same unknown normalization factor $Z_q$, so
$\hat G^{-1}\hat M\to(G/Z_q)^{-1}(M/Z_q)=G^{-1}M$: the normalization cancels exactly and no wave-function norm has to be computed. The projected columns then enter the same NES determinant and local-energy construction as in Eq.~\eqref{eq:local}, with an active block whose size is the number of active states and with the frozen columns entering as a constant reference block. Because the determinant is invariant under adding frozen components to an active column, the projection is enforced automatically and differentiating the determinant gives the projected tangent-space gradient; no explicit projection operator is ever formed. Physical energies and states are finally extracted from the generalized eigenvalue problem
\begin{equation}
M_\chi v=\lambda S_\chi v,
\qquad
(S_\chi)_{mn}=\braket{\chi_m}{\chi_n},
\qquad
(M_\chi)_{mn}=\bramat{\chi_m}{H}{\chi_n},
\label{eq:extraction}
\end{equation}
estimated on consistently weighted samples, so that the same normalization cancellation applies. At no point does the algorithm construct a $2^N$-dimensional vector, a dense Hamiltonian, or a dense projection operator.

The contributions of this work are as follows.
\begin{enumerate}
\item We formulate matrix-free deflation for frozen-subspace NES-VMC. The frozen states are reused through pointwise evaluations and a small rotation, the projection coefficients are estimated by Monte Carlo Gram matrices with exact normalization cancellation, and the projected active columns are optimized inside the NES determinant. The formulation is independent of the Hilbert-space dimension and requires only the ability to evaluate the frozen and active networks and to apply the Hamiltonian through its local connections.
\item We give a complete training algorithm that keeps sampling, local-energy estimation, and gradients consistent. It combines an auxiliary sampler for the projection coefficients, exponential moving averages of the column scales with the corresponding transformation of the projection coefficients, importance weighting to avoid zero-variance freezing in the replicated MCMC, real-parameter Wirtinger gradients and stochastic reconfiguration, and matrix-free extraction of the physical states.
\item We analyze the accuracy of the approach and identify the correct energy estimator. The eigenvalues of the sample-averaged local-energy matrix, $\operatorname{eig}\langle E_L\rangle$, reproduce the trace of the loss but are not the subspace eigenvalues for $K\ge2$; they can be off by hundreds of mHa even with exact projection coefficients. The consistently weighted Gram generalized-eigenvalue estimator of Eq.~\eqref{eq:extraction} is the appropriate quantity for energies and states. We also use an oracle diagonalization of the projected Hamiltonian $QHQ$ to separate the intrinsic frozen-subspace bias from the optimization error.
\item We demonstrate the method on H$_2$/6-31G in a CAS(2e,4o) active space (eight spin orbitals and a 16-dimensional FCI space). Freezing four states and optimizing two additional states reaches $E_4$ and $E_5$ within $0.0127$ and $0.0243$ mHa, with FCI overlaps of $0.99995$ and $0.99751$ and energy variances below $1.5\times10^{-5}$ Ha$^2$. Freezing four states and optimizing four additional states reaches $E_4$--$E_7$ within $0.003$--$0.010$ mHa, and 200 optimization steps already suffice at the chemical-accuracy scale. The corresponding full $K=6$ run, by contrast, left the errors in $E_4$ and $E_5$ at $59.7$ and $493$ mHa after 300 steps. The $K=2$ active run requires $608$ s for 1500 steps and is $1.7\times$ faster than a determinant-only matrix-free variant that keeps the frozen states inside a larger joint determinant. In the matrix-free runs, dense Hamiltonians and exact FCI vectors are used only for diagnostics and final verification, never inside the training loop.
\item We document the practical limitations of the approach and the diagnostics needed for reliable use: the frozen-subspace lower bound, residual overlap with the frozen span, condition numbers of the projected Gram matrices, drift of the unconstrained frozen components, and the coverage and effective sample size of the auxiliary distribution. These diagnostics distinguish a genuine optimization failure from an estimator or monitoring artifact.
\end{enumerate}

The remainder of the paper is organized as follows. Section~2 reviews NES-VMC and the matrix-free representation of the Hamiltonian. Section~3 develops the frozen-subspace deflation formalism, the Monte Carlo Gram estimators, and the projected gradient. Section~4 describes the training algorithm and its numerical safeguards. Section~5 presents the H$_2$/6-31G results, comparisons with full $K=6$ optimization and with determinant-only matrix-free deflation, and the estimator-bias analysis. Section~6 discusses the scope of the method, its limitations for larger active spaces, and the outlook for applications to LiH and other open-shell systems. Section~7 concludes.

\bibliographystyle{unsrtnat}
\bibliography{references}

\end{document}
